\documentclass[11pt]{article}
\newcommand{\documenttype}{Document}
\usepackage[a4paper,margin=0.85in]{geometry}
\usepackage[T1]{fontenc}
\usepackage[utf8]{inputenc}
\usepackage{lmodern,amsmath,amssymb,microtype,xcolor,enumitem,tabularx,titlesec,fancyhdr,hyperref}
\definecolor{accent}{HTML}{315B8A}
\definecolor{link}{HTML}{5B4F7A}
\hypersetup{colorlinks=true,urlcolor=link,linkcolor=link,citecolor=link,pdfauthor={Phurailatpam Hemantakumar}}
\setlength{\parindent}{0pt}\setlength{\parskip}{5pt}\setlength{\headheight}{14pt}\setlength{\emergencystretch}{2em}
\titleformat{\section}{\large\bfseries\color{accent}}{}{0pt}{}
\titlespacing*{\section}{0pt}{10pt}{5pt}
\setlist[itemize]{leftmargin=16pt,itemsep=2pt,topsep=3pt,parsep=0pt}
\pagestyle{fancy}\fancyhf{}\fancyhead[L]{\small Phurailatpam Hemantakumar}\fancyhead[R]{\small\documenttype}\fancyfoot[R]{\small\thepage}\renewcommand{\headrulewidth}{0.3pt}
\newcommand{\firstheading}[2]{\thispagestyle{plain}{\LARGE\bfseries\color{accent}#1}\par\vspace{3pt}{#2}\par\vspace{4pt}{\color{accent}\hrule}\vspace{8pt}}
\newcommand{\dated}[2]{\par\noindent\begin{tabularx}{\linewidth}{@{}X r@{}}#1&#2\end{tabularx}\par}
\newcommand{\contact}{\href{mailto:hemantaph@link.cuhk.edu.hk}{hemantaph@link.cuhk.edu.hk}\quad|\quad +852 9867 8178\quad|\quad\href{https://hemantaph.com}{hemantaph.com}\\Department of Physics, The Chinese University of Hong Kong}
\renewcommand{\documenttype}{Research statement}
\hypersetup{pdftitle={Research interests and plan | Monash OzGrav}}
\setlength{\parskip}{5pt}
\begin{document}
\firstheading{Research interests and proposed directions}{Phurailatpam Hemantakumar\quad|\quad Monash OzGrav Research Fellow (698908)}
\textbf{Research focus.} Gravitational-wave source populations, Bayesian inference and detector selection; strong gravitational lensing; eccentric compact-binary waveforms; efficient methods for large-scale population simulation.

My research connects astrophysical source models to signals observed by gravitational-wave (GW) detectors. I develop analytical and computational models, use them in Monte Carlo populations and Bayesian analyses, and test the reliability of resulting astrophysical conclusions. At Monash's OzGrav node, I would pursue two connected directions: population inference and lensing statistics, and waveform-based inference of compact-binary dynamics. These projects use my existing tools while creating natural opportunities for collaboration on GW source astrophysics, search sensitivity, and tests of gravity.

\section*{1. Previous work and capabilities}
\textbf{Lensing statistics and selection effects.} I lead the development of \href{https://ler.hemantaph.com/}{\texttt{ler}} (JOSS, 2026), an open-source framework for simulating lensed and unlensed compact-binary populations and predicting detectable rates. Its modelling includes source and lens distributions, elliptical power-law lenses with external shear, image properties, numerical cross-sections, optical depths, and GW selection effects. I have implemented propagation of source/lens hyperparameter and finite-sampling uncertainties; validation and publication of this extension are ongoing. The results can be used in lensing candidate analyses, selection functions and population likelihoods.

\textbf{Detection probability.} I lead development of \href{https://gwsnr.hemantaph.com/}{\texttt{gwsnr}}, which provides efficient matched-filter SNR calculations and approximate detection probabilities through waveform interpolation, accelerated inner products, and neural-network methods. Its computational efficiency enables large injection populations, but reliable astrophysical inference requires validation against actual search-recovery behaviour, particularly near threshold.

\textbf{Eccentric inspirals.} With collaborators I developed \texttt{TaylorF2Ecck}, a frequency-domain post-Newtonian model for nonspinning, low-eccentricity systems, and performed Bayesian analyses of GW170817 and GW190425 (Phys. Rev. D, 2026). Related work on GW190814 investigated degeneracy between eccentricity and line-of-sight acceleration (ApJ, 2026). These projects provide experience testing waveform approximations through consistency checks, injections, and event-level parameter estimation.

\newpage
\section*{2. Research directions at OzGrav}
\textbf{A. Population-constrained compact-binary rates and lensing forecasts.} I would start from star-formation histories constrained by electromagnetic observations and map them into compact-binary merger-rate densities through a small number of phenomenological parameters describing formation efficiency and delay times. Rather than fixing a single merger-rate curve, I would propagate the star-formation posterior into hierarchical inference with LVK compact-binary observations, consistently including detector selection. The central questions are which features of merger-rate evolution can be constrained by the GW catalogues, how much the answer depends on formation assumptions, and what these uncertainties imply for the observed strongly lensed population. \texttt{ler} would translate these inferred rates into lensing optical depths, image multiplicities, magnifications, and detectable rates, while propagating lens-population uncertainties. This framework could be tested with simulated catalogues and then applied to forthcoming LVK data releases.

\textbf{B. Calibrating selection functions for population inference.} A second objective is to test fast SNR-based detectability estimates against search-pipeline injection recoveries and their false-alarm-rate thresholds. I would quantify deviations across binary masses, distances, orientations and detector networks, then develop a calibration or interpolation scheme with propagated errors. This connects \texttt{gwsnr} to realistic search selection functions and helps distinguish astrophysical effects from sensitivity-related population distortions. It will also be valuable for lensed-event rate forecasts where magnification moves otherwise inaccessible events across the detection boundary.

\textbf{C. Eccentricity, acceleration and waveform systematics.} I would continue validation and development of low-eccentricity models, including higher-order eccentric effects and extensions incorporating spin where available. Injection studies and Bayesian model comparisons would test how waveform inaccuracies or external line-of-sight acceleration can mimic eccentricity, bias inferred source properties, or affect tests of gravity. The aim is to identify which additional waveform physics is required by present and future detector sensitivities, with transparent bounds on the validity of each approximation.

\section*{3. Expected outcomes and collaborations}
The proposed work is designed to produce reusable population likelihoods and selection functions, validated lensing-rate forecasts, tested waveform methods, and refereed analyses that can be incorporated into LVK projects. I would continue my existing LVK/GRASP responsibilities for \texttt{ler} and \texttt{gwsnr} while contributing to shared research at OzGrav. Monash's strengths in GW population modelling, Bayesian analysis, and compact-binary astrophysics make it an excellent environment to combine these threads. My experience with scientific Python, \texttt{bilby}, \texttt{LALSuite}, and accelerated computing would also support practical collaborations, research mentoring, and reproducible analysis workflows.
\end{document}
